\documentclass[11pt]{article}
\usepackage[utf8]{inputenc}
\usepackage{amsmath, amssymb, amsthm}
\usepackage{geometry}
\geometry{a4paper, margin=1in}
\usepackage{graphicx}
\usepackage{hyperref}
\usepackage{cite}
\usepackage{setspace}
\onehalfspacing

\title{Literature Review: Diffusion Models for Large Language Models (v3 -- DiDAL revision + corpus integration)}
\author{Your Name}
\date{\today}

\begin{document}

\maketitle

\begin{abstract}
Diffusion models now rival autoregressive models in continuous domains such as image, video, and audio synthesis, and have recently been extended to discrete data---most notably language---giving rise to diffusion language models, which the literature refers to interchangeably as Diffusion LLMs (DLLMs) and discrete diffusion language models (dLLMs). Unlike the dominant autoregressive (AR) paradigm, which generates tokens sequentially in a left-to-right fashion, diffusion-based language models generate text through an iterative denoising process that operates on entire sequences in parallel \cite{yu2025discrete, tseng2025diffusion}. This survey traces their evolution from continuous-space foundations to discrete formulations, categorizes representative models and techniques by diffusion space, generation strategy, and architecture, and examines applications, challenges, and future directions. We organize the field around the design choices that govern its quality--latency trade-off: the noise formulation, the generation strategy, and the post-training recipe.
\end{abstract}

\section{Introduction}

Autoregressive (AR) transformers remain the dominant paradigm for large language models, a position earned through scale and mature training infrastructure rather than an inherent advantage of sequential decoding \cite{yu2025discrete}. AR models generate text token by token and achieve remarkable performance across a wide range of tasks, but they possess intrinsic limitations: left-to-right decoding restricts parallelization during inference \cite{yu2025discrete}, precise structural control (e.g., length or format constraints) is difficult to enforce during sampling, and dynamic, task-aware perception requires costly chain-of-thought or multi-round processing \cite{yu2025discrete}.

Diffusion models---originally developed for continuous domains such as image generation \cite{ho2020denoising}---offer a compelling alternative. By framing generation as an iterative denoising process, diffusion models enable parallel token generation, bidirectional context modeling, and fine-grained controllability \cite{li2025surveydlm}. The application of diffusion to language has progressed from small-scale experiments to commercial-scale deployments \cite{khanna2025mercury}, with commercial deployments reporting roughly an order-of-magnitude latency reduction over speed-optimized AR baselines under batch decoding \cite{yu2025discrete, wang2025d2f}.

This literature review synthesizes the current state of diffusion models for LLMs, organized as follows: Section 2 covers the mathematical and conceptual foundations. Section 3 traces the evolution from continuous to discrete formulations. Section 4 categorizes representative models and techniques. Section 5 examines applications and empirical results. Section 6 discusses challenges and open problems. Section 7 outlines future directions.

\section{Foundations of Diffusion Models}

\subsection{The Diffusion Framework}

Diffusion models are generative models that learn to reverse a gradual noising process \cite{ho2020denoising}. The framework consists of two core processes:

\textbf{Forward Process (Noising):} Starting from clean data $\mathbf{x}_0$, a Markov chain gradually adds noise over $T$ timesteps, producing increasingly corrupted versions $\mathbf{x}_1, \mathbf{x}_2, \ldots, \mathbf{x}_T$ until the data approximates pure noise \cite{tseng2025diffusion}.

\textbf{Reverse Process (Denoising):} A neural network is trained to invert this process, learning to predict and remove noise step by step, ultimately generating new samples from random noise \cite{ho2020denoising}.

Three foundational works shaped the continuous diffusion framework. The Denoising Diffusion Probabilistic Model (DDPM) established the core variational framework: training reduces to a simple weighted mean-squared-error objective on noise prediction, which made large-scale training practical for the first time \cite{ho2020denoising}. Denoising Diffusion Implicit Models (DDIM) generalized sampling to non-Markovian trajectories, enabling deterministic sampling with orders-of-magnitude fewer steps and thus making iterative refinement computationally viable at inference time \cite{song2021ddim}. Score-based generative modeling through stochastic differential equations (SDEs) then unified DDPM and score matching under a continuous-time formulation, providing the theoretical lens through which modern samplers---including the discrete formulations discussed below---are analyzed \cite{song2021scoresde}. For discrete state spaces, the analogous continuous-time object is a continuous-time Markov chain (CTMC): Campbell et al.\ derived a continuous-time evidence lower bound for discrete denoising models with tau-leaping and predictor-corrector samplers \cite{campbell2022ctmc}, and Sun et al.\ developed the corresponding score-based continuous-time formulation \cite{sun2022ctmc}. These CTMC frameworks underpin the continuous-time objectives of modern masked diffusion LLMs.

\subsection{The Challenge of Discrete Data}

A fundamental challenge arises when applying diffusion to language: text is inherently discrete (composed of tokens), whereas diffusion models were originally designed for continuous data (e.g., image pixels) \cite{shen2026codar}. In continuous domains, Gaussian noise can be naturally added and removed; in discrete domains, adding noise to a token renders it invalid \cite{jin2025discreteness}.

This challenge has led to two distinct approaches \cite{jin2025discreteness}:

\begin{enumerate}
    \item \textbf{Continuous Diffusion in Embedding Space:} Mapping discrete tokens to continuous embeddings and applying standard diffusion in this space \cite{shen2026codar}.
    \item \textbf{Discrete Diffusion over Tokens:} Defining noise processes directly over discrete token spaces using transition matrices \cite{tseng2025diffusion}.
\end{enumerate}

To handle discrete data directly, Austin et al. proposed Structured Denoising Diffusion Models in Discrete State-Spaces (D3PM), which generalizes the forward process to a family of discrete transition matrices---including uniform and absorbing-state (masking) kernels \cite{austin2021d3pm}. Concurrently, Hoogeboom et al.\ introduced multinomial diffusion with argmax flows for categorical data, an earlier and closely related formulation \cite{hoogeboom2021multinomial}.

\section{Evolution of Diffusion Language Models}

The evolution of diffusion models for language can be traced through several distinct phases \cite{tseng2025diffusion}.

\subsection{Early Foundations (2015-2020)}

The thermodynamic origins of diffusion models were established through DDPM \cite{ho2020denoising} and DDIM \cite{song2021ddim}, which introduced the core principles of Markovian noising and denoising on continuous data, later unified in continuous time by the score-SDE formulation \cite{song2021scoresde}. These foundational works demonstrated that iterative refinement from noise could produce high-quality samples, though they were initially applied exclusively to continuous domains like images.

\subsection{First Text-Specific Adaptations (2021-2022)}

The first systematic efforts to adapt diffusion for text emerged in this period. Zou et al. provided the first systematic survey of diffusion for non-autoregressive text generation, defining both discrete (mask-based) and continuous (Gaussian embedding) formulations \cite{zou2023survey}.

Key pioneering works include:

\textbf{D3PM} (Austin et al., 2021): This work proposed a diffusion-like generative model for discrete data, defining a family of discrete damage processes that generalize the diffusion framework beyond Gaussian noise \cite{austin2021d3pm}. It demonstrated that absorbing-state discrete diffusion can produce coherent categorical samples without timestep inputs \cite{austin2021d3pm}, establishing the technical foundation on which most modern dLLMs rest.

\textbf{Diffusion-LM} (Li et al., 2022): This work developed a non-autoregressive language model based on continuous diffusions, embedding Gaussian noise into token representations \cite{li2022diffusion, tseng2025diffusion}. Diffusion-LM demonstrated successful control for six challenging fine-grained control tasks---including syntactic structure, length, and lexical constraints---significantly outperforming prior plug-in methods in controllable text generation \cite{li2022diffusion}.

\textbf{DiffusionBERT} (Wang et al., 2022): This work presented a generative masked language model based on discrete diffusion models, leveraging BERT as its backbone \cite{wang2023diffusionbert}. By recognizing that diffusion models and pre-trained language models share a denoising objective, DiffusionBERT achieved significant improvements over existing diffusion models for text \cite{wang2023diffusionbert}.

\textbf{SSD-LM} (Han et al., 2022): This work introduced a semi-autoregressive diffusion-based language model that iteratively generates blocks of text, allowing for flexible output length while enabling local bidirectional context updates within each block \cite{han2023ssdlm}. Its block-wise design is architecturally significant: rather than treating parallelism and sequential structure as mutually exclusive, SSD-LM shows they can be combined. The work also demonstrated that with simplex projection, the diffusion process can be conducted directly in the natural vocabulary space \cite{han2023ssdlm}.

\textbf{CDCD} (Dieleman et al., 2022): This framework modeled categorical data with diffusion processes that are continuous both in time and input space \cite{dieleman2022cdcd}, demonstrating efficacy on several language modeling tasks and providing score interpolation techniques for improved likelihood estimation \cite{dieleman2022cdcd}.

\subsection{The Shift to Discrete Diffusion (2023-2024)}

Continuous approaches faced a persistent difficulty: mapping continuous embeddings back to discrete tokens without quality degradation, since small perturbations in embedding space can decode to unrelated tokens. Recent analyses suggest this gap was driven at least as much by unexamined design choices in early continuous systems as by any fundamental property of discreteness \cite{shen2026codar}, but empirically research shifted decisively toward discrete formulations during this period \cite{jin2025discreteness}.

\textbf{Masked Diffusion Language Models (MDLM)} (Sahoo et al., 2024): This work introduced a masked discrete diffusion model featuring a novel substitution-based parameterization that simplifies the absorbing-state diffusion loss to a Rao--Blackwellized mixture of classical masked language modeling cross-entropies \cite{sahoo2024simple}. MDLM demonstrated that masked diffusion could close much of the perplexity gap to AR models at small scales \cite{sahoo2024simple}.

\textbf{SEDD} (Lou et al., 2023): Rather than parameterizing a vocabulary-sized transition matrix, SEDD estimates the ratios of the data distribution across noise scales with a score-entropy objective \cite{lou2023sedd}. This ratio-based view is the direct ancestor of the score parameterizations used in subsequent scaled masked diffusion models, and SEDD showed competitive likelihoods against GPT-2 at its scale \cite{lou2023sedd}.

\textbf{GIDD} (von R\"utte et al., 2025): Generalized interpolating discrete diffusion unifies masked and uniform noising in a single forward-process family, including hybrid schedules with self-correction capability \cite{gulrajani2025gidd}. In the reported small-model studies, mask-only models generally performed better on likelihood and downstream tasks; whether uniform components help token-constrained scaling remains an open hypothesis \cite{gulrajani2025gidd}.

\subsection{Industrial-Scale Diffusion LLMs (2025-Present)}

The most recent phase has witnessed the scaling of diffusion language models to commercial scale, with performance rivaling or matching autoregressive counterparts \cite{yu2025discrete}.

\textbf{Mercury} (Inception Labs, 2025): The first commercially available diffusion-based LLM, Mercury achieves throughputs of 1,109 tokens/sec (Mini) and 737 tokens/sec (Small) on NVIDIA H100 GPUs---up to $10\times$ faster than speed-optimized frontier models---while maintaining comparable quality on coding benchmarks \cite{khanna2025mercury}. Mercury Coder Mini achieves an 88\% score on HumanEval and ranks second on Copilot Arena with just 25ms latency \cite{khanna2025mercury}.

\textbf{LLaDA} (Nie et al., 2025): The first diffusion-based language model to achieve performance comparable to autoregressive models at the 8B scale, trained entirely from scratch with fully bidirectional attention \cite{nie2025large}. LLaDA2.0 later scaled to 100B parameters, representing the first diffusion model at this scale \cite{nie2025large}. Complementary scaling work adapted existing AR checkpoints to the diffusion objective: DiffuLLaMA reached 7B by adaptation \cite{gong2024diffullama}, and the masked-diffusion MDM study reported a loss scaling rate similar to AR models while quantifying a roughly $16\times$ compute gap at matched validation loss \cite{nie2024mdm}. Earlier likelihood-based diffusion approaches faced a substantially larger ($\sim$64$\times$) matched-likelihood disadvantage \cite{gulrajani2023plaid}, so these gaps reflect different objectives and estimators rather than a single constant penalty. Compute-optimal training studies further find diffusion-specific tradeoffs, including higher optimal data budgets than AR models in single-epoch regimes \cite{ni2025quokka}.

\subsection{Diffusion-Native Post-Training (2025-Present)}

Adapting preference optimization and reinforcement learning to diffusion decoders has become its own research thread. SEPO applies policy gradients to discrete diffusion with self-normalized importance sampling and clipped ratios \cite{zekri2025sepo}; d1 adapts GRPO to masked dLLMs through one-step unmasking with verifiable rewards \cite{zhao2025d1}; VRPO reduces the variance of ELBO-based DPO estimators through Monte Carlo allocation and antithetic sampling, yielding LLaDA 1.5 \cite{zhu2025llada15}; and coupled-GRPO uses complementary mask noise for paired completions in code models such as DiffuCoder \cite{gong2025diffucoder}. Reported gains are promising but rest on limited data, task, and model coverage.

\section{Categorization of Diffusion Language Models}

\subsection{By Diffusion Space}

\textbf{Continuous-Space Models:} These models map discrete tokens to continuous embeddings and apply standard Gaussian diffusion \cite{shen2026codar}. Examples include Diffusion-LM \cite{li2022diffusion}, SSD-LM \cite{han2023ssdlm}, and CDCD \cite{dieleman2022cdcd}. While these models benefit from the well-understood mathematics of continuous diffusion, they face challenges in maintaining semantic consistency when projecting back to discrete tokens \cite{shen2026codar}.

\textbf{Discrete-Space Models:} These models operate directly on token spaces using discrete noise processes such as absorbing-state (masking) or uniform transitions \cite{jin2025discreteness}. Examples include D3PM \cite{austin2021d3pm}, MDLM \cite{sahoo2024simple}, DiffusionBERT \cite{wang2023diffusionbert}, Mercury \cite{khanna2025mercury}, and LLaDA \cite{nie2025large}. Discrete models have largely superseded continuous approaches in recent work \cite{jin2025discreteness}.

\subsection{By Generation Strategy}

\textbf{Fully Parallel Models:} These models generate all tokens simultaneously through iterative refinement. Examples include MDLM \cite{sahoo2024simple} and LLaDA \cite{nie2025large}.

\textbf{Semi-Autoregressive Models:} These models generate blocks of text iteratively, combining the parallelism of diffusion with the sequential structure of autoregression. SSD-LM exemplifies this approach \cite{han2023ssdlm}.

\subsection{By Architecture}

\textbf{Transformer-Based:} Most diffusion language models employ Transformer architectures, as the choice of architecture is orthogonal to the diffusion paradigm \cite{khanna2025mercury}. Mercury, LLaDA, and DiffusionBERT all use Transformer backbones.

\textbf{Hybrid Architectures:} Some works explore combining diffusion with other paradigms, such as integrating LLMs with diffusion models for multimodal generation \cite{benjdira2025integrating}.

\section{Applications and Empirical Results}

\subsection{Code Generation}

Coding represents a particularly latency-sensitive domain where diffusion models have demonstrated exceptional value \cite{khanna2025mercury}. Mercury Coder achieves:

\begin{itemize}
    \item \textbf{HumanEval:} 88\% (Mini) and 90\% (Small) pass@1 \cite{khanna2025mercury}
    \item \textbf{MultiPL-E:} Competitive performance across C++, Java, JavaScript, PHP, Bash, and TypeScript \cite{khanna2025mercury}
    \item \textbf{Fill-in-the-Middle:} State-of-the-art performance of 93.1\% (Small) on single-line FIM \cite{khanna2025mercury}
    \item \textbf{Speed:} 1,109 tokens/sec (Mini) and 737 tokens/sec (Small) on H100 GPUs \cite{khanna2025mercury}
\end{itemize}

These results demonstrate that diffusion models can match or exceed AR models in code generation while achieving order-of-magnitude speed improvements. The throughput figures should be read in context: they are aggregate batch throughputs on H100 GPUs, and independent analyses attribute part of the speedup to the reduced number of denoising steps (Mercury uses a 10-step schedule) rather than to parallel decoding alone \cite{wang2025d2f, fu2025efficientdlm}.

\subsection{Controllable Text Generation}

One of the earliest and most compelling applications of diffusion to language was controllable generation. Diffusion-LM demonstrated successful control across six fine-grained tasks, significantly outperforming prior work \cite{li2022diffusion}. The bidirectional nature of diffusion enables constraints to be applied throughout the generation process rather than only left-to-right \cite{li2025surveydlm}.

\subsection{Multimodal Generation}

The integration of LLMs and diffusion models has emerged as a key research frontier for advancing multimodal generation, reasoning, and cross-domain understanding \cite{benjdira2025integrating}. This integration enables:

\begin{itemize}
    \item \textbf{Text-to-Image Synthesis:} Leveraging LLMs for semantic understanding and diffusion models for high-quality generation \cite{benjdira2025integrating}
    \item \textbf{Text-to-Video Synthesis:} Extending diffusion to temporal domains with LLM guidance \cite{benjdira2025integrating}
    \item \textbf{Unified Multimodal Models:} Combining understanding and generation capabilities in a single framework \cite{khanna2025mercury}
\end{itemize}

\subsection{Human Evaluation}

Mercury Coder Mini, evaluated on Copilot Arena, ranks tied for second place in code completion quality, surpassing speed-optimized models like GPT-4o Mini and Gemini-1.5-Flash, and even larger models like GPT-4o \cite{khanna2025mercury}. With an average latency of just 25ms, it is approximately $4\times$ faster than GPT-4o Mini, which averages around 100ms in the same evaluation \cite{khanna2025mercury}.

\section{Challenges (Open Problems) vs.\ Future Directions}

Because the remaining two analytical sections address related but distinct concerns, we state their partition explicitly: this section catalogs \emph{unsolved problems with evidence}, while Section 7 outlines \emph{proposed research programs}. Each challenge below maps to one or more directions that follow.

\subsection{Challenges and Open Problems}

\subsection{Inference Efficiency}

While diffusion models offer theoretical parallelism, practical efficiency depends on the number of denoising steps. Early diffusion models required hundreds of steps; training-free acceleration techniques such as KV caching and parallel decoding (Fast-dLLM), as well as learned step-distillation samplers, have since reduced this to 10-20 steps with minimal quality loss \cite{wu2025fastdllm, fu2025efficientdlm}. The trade-off between generation quality and inference speed remains an active area of research \cite{li2025surveydlm}.

\subsection{Long-Sequence Handling}

Diffusion models face challenges in handling long sequences due to the computational cost of processing entire sequences in each denoising step \cite{li2025surveydlm}. Context windows currently range from 32K (Mercury) \cite{khanna2025mercury} to 128K with positional-extension techniques such as RoPE scaling, compared to 1M+ for some AR models.

\subsection{Scaling Laws}

The scaling properties of diffusion language models are less well understood than those of autoregressive models \cite{khanna2025mercury}. While larger diffusion models consistently outperform smaller ones, the relationship between model size, data, and performance requires further investigation \cite{yu2025discrete}.

\subsection{Alignment and Fine-Tuning}

Adapting diffusion models to follow instructions through RLHF or DPO presents unique challenges, as the denoising objective differs fundamentally from the autoregressive next-token prediction loss \cite{khanna2025mercury}. Techniques developed for AR models may not transfer directly.

\subsection{Evaluation Protocols}

Standardized benchmarks and evaluation protocols for diffusion language models remain underdeveloped \cite{tseng2025diffusion}. Many existing benchmarks were designed for AR models and may not capture the unique strengths and weaknesses of diffusion-based generation \cite{yu2025discrete}.

\section{Future Directions}

\subsection{Unified Autoregressive-Diffusion Models}

Recent work has explored unifying AR and diffusion paradigms through hyperschedules and hybrid architectures \cite{tseng2025diffusion}. Such unification could enable models to leverage the strengths of both approaches: the deep reasoning of AR and the speed and controllability of diffusion. This direction directly addresses the scaling-law and alignment challenges of Section 6, since hybrid schedules can retain AR-style likelihood training while inheriting diffusion-style parallel decoding.

\subsection{Efficient Inference Techniques}

Advances in caching mechanisms \cite{wu2025fastdllm}, adaptive correction sampling \cite{tseng2025diffusion}, and decoding parallelism \cite{li2025surveydlm} promise to further reduce inference latency. The development of specialized hardware and inference engines for diffusion models represents another important direction \cite{khanna2025mercury}.

\subsection{Scaling to Frontier Performance}

While diffusion models have demonstrated comparable performance to speed-optimized AR models, matching frontier models (e.g., GPT-4o, Claude 3.5 Sonnet) on complex reasoning benchmarks remains an open challenge \cite{khanna2025mercury}. Scaling diffusion models to hundreds of billions of parameters, as demonstrated by LLaDA2.0 \cite{nie2025large}, may help close this gap.

\subsection{Agentic Applications}

The speed advantages of diffusion models make them particularly attractive for agentic workloads requiring multiple reasoning and execution cycles \cite{khanna2025mercury}. However, memory and context limitations must be addressed for complex, multi-step agentic tasks.

\subsection{Multimodal Unification}

The integration of diffusion models and LLMs for unified multimodal understanding and generation represents a promising frontier \cite{benjdira2025integrating, khanna2025mercury}. Such unified models could serve as the foundation for next-generation AI systems capable of seamless interaction across modalities.

\section{Conclusion}

Diffusion models for large language models have evolved from small-scale experiments to commercial-scale deployments in just a few years. The shift from continuous to discrete formulations---driven by the token-projection difficulties of embedding-space approaches, though recent evidence suggests the gap was narrower than commonly assumed \cite{shen2026codar}---coupled with advances in training and inference efficiency, has enabled diffusion-based LLMs to achieve performance comparable to autoregressive models while delivering order-of-magnitude speed improvements under batch decoding \cite{yu2025discrete, wang2025d2f}.

Key milestones include the pioneering Diffusion-LM \cite{li2022diffusion}, the discrete diffusion frameworks of D3PM \cite{austin2021d3pm} and MDLM \cite{sahoo2024simple}, the industrial-scale deployment of Mercury \cite{khanna2025mercury}, and the 100B-scale LLaDA2.0 \cite{nie2025large}. These developments position diffusion models as a compelling alternative---and complement---to the dominant autoregressive paradigm \cite{yu2025discrete}.

However, significant challenges remain, including inference efficiency optimization, long-sequence handling, scaling law characterization, alignment techniques, and standardized evaluation protocols \cite{tseng2025diffusion, li2025surveydlm}. Addressing these challenges will be essential for realizing the full potential of diffusion-based language models in real-world applications.

The future of language modeling likely lies not in a single paradigm but in hybrid approaches that leverage the complementary strengths of autoregressive and diffusion models \cite{benjdira2025integrating}---combining the deep reasoning of AR with the speed, parallelism, and controllability of diffusion.

\bibliographystyle{plain}
\begin{thebibliography}{99}

\bibitem{ho2020denoising}
J. Ho, A. Jain, and P. Abbeel, ``Denoising diffusion probabilistic models,'' \textit{Advances in Neural Information Processing Systems}, vol. 33, pp. 6840--6851, 2020.

\bibitem{li2022diffusion}
X. L. Li, J. Thickstun, I. Gulrajani, P. Liang, and T. B. Hashimoto, ``Diffusion-LM improves controllable text generation,'' \textit{arXiv preprint arXiv:2205.14217}, 2022.

\bibitem{tseng2025diffusion}
C.-Y. Tseng, D. Zhang, Z. Bi, and J. Song, ``Diffusion-based large language models survey,'' \textit{TechRxiv}, 2025.

\bibitem{yu2025discrete}
R. Yu, Q. Li, and X. Wang, ``Discrete diffusion in large language and multimodal models: A survey,'' \textit{arXiv preprint arXiv:2506.13759}, 2025.

\bibitem{li2025surveydlm}
T. Li, M. Chen, B. Guo, and Z. Shen, ``A survey on diffusion language models,'' \textit{arXiv preprint arXiv:2508.10875}, 2025.

\bibitem{benjdira2025integrating}
A. Benjdira and A. M. Ali, ``Integrating large language models and diffusion models in generative AI tasks: Progress, challenges, and future directions,'' \textit{TechRxiv}, 2025.

\bibitem{zou2023survey}
H. Zou, Z. M. Kim, and D. Kang, ``A survey of diffusion models in natural language processing,'' \textit{arXiv preprint arXiv:2305.14671}, 2023.

\bibitem{sahoo2024simple}
S. Sahoo et al., ``Simple and effective masked diffusion language models,'' \textit{NeurIPS}, 2024.

\bibitem{wang2023diffusionbert}
K. Wang et al., ``DiffusionBERT: Improving generative masked language models with diffusion models,'' \textit{ACL}, 2023.

\bibitem{han2023ssdlm}
X. Han et al., ``SSD-LM: Semi-autoregressive simplex-based diffusion language model for text generation and modular control,'' \textit{ACL}, 2023.

\bibitem{dieleman2022cdcd}
S. Dieleman et al., ``Continuous diffusion for categorical data,'' \textit{arXiv preprint}, 2022.

\bibitem{nie2025large}
Y. Nie et al., ``Large language diffusion models (LLaDA),'' \textit{NeurIPS}, 2025.

\bibitem{khanna2025mercury}
S. Khanna et al., ``Mercury: Ultra-fast language models based on diffusion,'' \textit{Inception Labs Technical Report}, 2025.

\bibitem{song2021ddim}
J. Song, C. Meng, and S. Ermon, ``Denoising diffusion implicit models,'' \textit{arXiv preprint arXiv:2010.02502}, 2021.

\bibitem{austin2021d3pm}
J. Austin, D. D. Johnson, J. Ho, D. Tarlow, and R. van den Berg, ``Structured denoising diffusion models in discrete state-spaces,'' \textit{arXiv preprint arXiv:2107.03006}, 2021.

\bibitem{song2021scoresde}
Y. Song, J. Sohl-Dickstein, D. P. Kingma, A. Kumar, S. Ermon, and B. Poole, ``Score-based generative modeling through stochastic differential equations,'' \textit{arXiv preprint arXiv:2011.13456}, 2021.

\bibitem{hoogeboom2021multinomial}
E. Hoogeboom, D. Nielsen, P. Jaini, P. Forr\'e, and M. Welling, ``Argmax flows and multinomial diffusion: Learning categorical distributions,'' \textit{arXiv preprint arXiv:2102.05379}, 2021.

\bibitem{shen2026codar}
J. Shen, J. Zhao, Z. He, and Z. Lin, ``CoDAR: Continuous diffusion language models are more powerful than you think,'' \textit{arXiv preprint arXiv:2603.02547}, 2026.

\bibitem{jin2025discreteness}
Z. Jin, B. Wang, X. Lin, L. Bing, and A. Sun, ``On the role of discreteness in diffusion LLMs,'' \textit{arXiv preprint arXiv:2512.22630}, 2025.

\bibitem{wang2025d2f}
X. Wang, C. Xu, Y. Jin, J. Jin, H. Zhang, and Z. Deng, ``Diffusion LLMs can do faster-than-AR inference via discrete diffusion forcing,'' \textit{arXiv preprint arXiv:2508.09192}, 2025.

\bibitem{fu2025efficientdlm}
Y. Fu, L. Whalen, Z. Ye, X. Dong, S. Diao, and J. Liu, ``Efficient-DLM: From autoregressive to diffusion language models, and beyond in speed,'' \textit{arXiv preprint arXiv:2512.14067}, 2025.

\bibitem{wu2025fastdllm}
C. Wu, H. Zhang, and S. Xue, ``Fast-dLLM: Training-free acceleration of diffusion LLM by enabling KV cache and parallel decoding,'' \textit{arXiv preprint arXiv:2505.22618}, 2025.

\bibitem{campbell2022ctmc}
A. Campbell, J. Benton, V. De Bortoli, T. Rainforth, G. Deligiannidis, and A. Doucet, ``A continuous time framework for discrete denoising models,'' \textit{arXiv preprint arXiv:2205.14987}, 2022.

\bibitem{sun2022ctmc}
H. Sun, L. Yu, B. Dai, D. Schuurmans, and H. Dai, ``Score-based continuous-time discrete diffusion models,'' \textit{arXiv preprint arXiv:2211.16750}, 2022.

\bibitem{lou2023sedd}
A. Lou, C. Meng, and S. Ermon, ``Discrete diffusion modeling by estimating the ratios of the data distribution,'' \textit{arXiv preprint arXiv:2310.16834}, 2023.

\bibitem{gulrajani2025gidd}
D. von R\"utte, J. Fluri, Y. Ding, A. Orvieto, B. Sch\"olkopf, and T. Hofmann, ``Generalized interpolating discrete diffusion,'' \textit{arXiv preprint arXiv:2503.04482}, 2025.

\bibitem{gulrajani2023plaid}
I. Gulrajani and T. B. Hashimoto, ``Likelihood-based diffusion language models,'' \textit{arXiv preprint arXiv:2305.18619}, 2023.

\bibitem{nie2024mdm}
S. Nie, F. Zhu, Z. Gao, D. Du, T. Pang, Q. Liu, C. Li, and M. Lin, ``Scaling up masked diffusion models on text,'' \textit{arXiv preprint arXiv:2410.18514}, 2024.

\bibitem{gong2024diffullama}
S. Gong, S. Agarwal, Y. Zhang, J. Ye, L. Zheng, M. Li, C. An, P. Zhao, and T. Pfister, ``Scaling diffusion language models via adaptation from autoregressive models,'' \textit{arXiv preprint arXiv:2410.17891}, 2024.

\bibitem{ni2025quokka}
J. Ni, Q. Liu, C. Du, L. Dou, H. Yan, Z. Wang, and W. Lu, ``Training optimal large diffusion language models,'' \textit{arXiv preprint arXiv:2510.03280}, 2025.

\bibitem{zekri2025sepo}
O. Zekri and N. Boull\'e, ``Fine-tuning discrete diffusion models with policy gradient methods,'' \textit{arXiv preprint arXiv:2502.01384}, 2025.

\bibitem{zhao2025d1}
S. Zhao, D. Gupta, Q. Zheng, and A. Grover, ``d1: Scaling reasoning in diffusion large language models via reinforcement learning,'' \textit{arXiv preprint arXiv:2504.12216}, 2025.

\bibitem{zhu2025llada15}
F. Zhu, R. Wang, S. Nie, X. Zhang, C. Wu, J. Chen, Y. Liu, and M. Lin, ``LLaDA 1.5: Variance-reduced preference optimization for large language diffusion models,'' \textit{arXiv preprint arXiv:2505.19223}, 2025.

\bibitem{gong2025diffucoder}
S. Gong, R. Zhang, H. Zheng, J. Gu, N. Jaitly, L. Kong, and Y. Li, ``DiffuCoder: Understanding and improving masked diffusion models for code generation,'' \textit{arXiv preprint arXiv:2506.20639}, 2025.

\end{thebibliography}

\end{document}
